\documentclass[11pt]{article}

\usepackage[a4paper,margin=2.4cm]{geometry}
\usepackage{amsmath,amssymb}
\usepackage{array}
\usepackage{booktabs}
\usepackage{enumitem}
\usepackage[hidelinks]{hyperref}
\usepackage{siunitx}

\newcommand{\file}[1]{\texttt{#1}}
\newcommand{\Ox}[1]{{}^{#1}\mathrm{O}}
\newcommand{\Cl}[1]{{}^{#1}\mathrm{C}}

\title{Implementation plan for proton-induced\\
cross-section systematics in \file{ptcryspg4}}
\author{J.J. G\'omez-Cadenas}
\date{\today}

\begin{document}
\maketitle

\begin{abstract}
The positron-emitter source in \file{ptcryspg4} is generated by Geant4 proton
transport with the \file{QGSP\_BIC\_HP} reference physics list.  This note
describes an implementation that propagates uncertainties in the production
cross sections to the reconstructed distal activity edge.  The primary method
is channel-aware reweighting of emitter production while retaining the nominal
Geant4 transport.  Complete simulations with alternative Geant4 cascade models
provide a complementary hadronic-model envelope.  The plan covers the required
production metadata, cross-section inputs, weighting, acquisition budgets,
validation and production sequence.
\end{abstract}

\section{Objective}

For a cross-section model $m$, the quantity to be propagated is the displacement
of the reconstructed distal activity edge relative to the nominal prediction,
\begin{equation}
  b_{\mathrm{xs}}^{(m)} =
  R_{50}^{(m)}-R_{50}^{(\mathrm{nominal})}.
  \label{eq:bias}
\end{equation}
Here $R_{50}$ is obtained with the same acquisition, detector simulation,
reconstruction and edge estimator used in the range-verification study.  The
set of shifts $b_{\mathrm{xs}}^{(m)}$ defines the cross-section contribution to
the distal-edge systematic uncertainty.

The implementation separates two related effects:
\begin{enumerate}[nosep]
  \item \emph{cross-section-only variations}, obtained by reweighting isotope
        production while retaining the nominal particle transport;
  \item a broader \emph{hadronic-model envelope}, obtained from full Geant4 runs
        with alternative cascade models.
\end{enumerate}

\section{Current Stage-A implementation}

The executable \file{stageA\_transport/proton\_transport.cc} constructs
\file{QGSP\_BIC\_HP} directly.  At the energies of the treatment field, Geant4
samples proton inelastic interactions and the Binary Cascade generates the
reaction final state.  A residual nucleus matching one of the five registered
positron emitters is recognized in \file{TrackingAction.cc} by its nuclear
charge and mass $(Z,A)$.  Its positron is subsequently transported to its
annihilation point.

The current \file{emitters.csv} row contains the primary event identifier,
isotope identifier, production coordinates and annihilation coordinates.  The
following production information is not retained:
\begin{itemize}[nosep]
  \item identity and kinetic energy of the projectile at the interaction;
  \item target-nucleus $Z$ and $A$;
  \item creator process and Geant4 final-state model;
  \item an explicit production-channel identifier.
\end{itemize}
The existing frozen \file{uniform\_headep\_sobp\_1e8} source therefore does not
contain enough information for reaction-by-reaction reweighting.  A new Stage-A
production is required after the provenance fields have been implemented.

The nominal SOBP samples incident energies from approximately
\SIrange{81.3}{112.9}{MeV}.  Transport through the phantom spans the lower
energies that determine isotope production near the distal edge.  The
cross-section representation must therefore extend from the reaction threshold
through the maximum incident energy.

\section{Production provenance}

\subsection{Emitter record}

Each emitter is to carry the compact numerical fields in
Table~\ref{tab:provenance}.  Reaction labels are inferred from the projectile,
target and residual, avoiding repeated strings in a source containing about two
million rows.

\begin{table}[htbp]
  \centering
  \caption{Additional production information for each emitter.}
  \label{tab:provenance}
  \begin{tabular}{lll}
    \toprule
    Field & Suggested type & Meaning \\
    \midrule
    \file{projectile\_pdg}       & \file{int32}   & projectile identity \\
    \file{projectile\_E\_MeV}   & \file{float32} & pre-interaction kinetic energy \\
    \file{target\_Z}             & \file{int16}   & target nuclear charge \\
    \file{target\_A}             & \file{int16}   & target mass number \\
    \file{residual\_Z}           & \file{int16}   & residual nuclear charge \\
    \file{residual\_A}           & \file{int16}   & residual mass number \\
    \file{creator\_process}      & compact code   & producing Geant4 process \\
    \file{creator\_model\_id}   & \file{int32}   & Geant4 final-state model ID \\
    \bottomrule
  \end{tabular}
\end{table}

For example,
\begin{equation}
  \bigl(p,\Ox{16},\Ox{15}\bigr)
  \quad\Longrightarrow\quad
  \Ox{16}(p,pn)\Ox{15}.
\end{equation}
The channels required initially are
\begin{align}
  \Ox{16}(p,pn)\Ox{15},\qquad
  \Cl{12}(p,pn)\Cl{11},\qquad
  \Ox{16}(p,3p3n)\Cl{11}.
  \label{eq:channels}
\end{align}

\subsection{Capture inside Geant4}

The production record can be assembled during the step that creates the
residual nucleus:
\begin{enumerate}[nosep]
  \item \file{SteppingAction} identifies a hadronic step that produces a listed
        residual among the secondaries.
  \item The pre-step track supplies the projectile identity and kinetic energy.
  \item \file{G4HadronicProcess::GetTargetNucleus()} supplies the sampled target
        $Z$ and $A$.
  \item The residual track supplies its track identifier, vertex, isotope and
        creator-model ID.
  \item \file{EventAction} retains the record, indexed by residual track ID.
  \item \file{TrackingAction} associates the later positron with the record and
        appends its annihilation point.
\end{enumerate}
This also makes the residual-ion vertex the stored production coordinate.  A
validation counter will report the fraction of emitter rows with complete
production provenance.

\section{Selectable Geant4 physics}

The physics list will become a run-time parameter, with
\file{QGSP\_BIC\_HP} retained as the default.  The Geant4
\file{G4PhysListFactory} can construct the requested reference list before
initialization.  A representative invocation is
\begin{verbatim}
./proton_transport --physics-list QGSP_INCLXX_HP \
    uniform_headep_sobp.mac
\end{verbatim}

The selected name will be passed to \file{RunAction}, recorded in
\file{run\_meta.csv}, and included in the run tag.  Including it in the tag
prevents two otherwise identical runs from sharing an output directory.

The initial model ensemble is
\begin{itemize}[nosep]
  \item \file{QGSP\_BIC\_HP}, the nominal Binary Cascade calculation;
  \item \file{QGSP\_INCLXX\_HP}, an INCL++ cascade calculation;
  \item \file{QGSP\_BERT\_HP}, a Bertini cascade calculation.
\end{itemize}
At the treatment energies, the cascade choice controls the residual-nucleus
final state.  These complete runs vary isotope production and the associated
hadronic transport together, and their spread is reported as a hadronic-model
envelope.

\section{Cross-section representation}

The input curves will be stored as versioned CSV data with the schema
\begin{verbatim}
model,projectile_pdg,target_Z,target_A,residual_Z,residual_A,
energy_MeV,sigma_mb,sigma_uncertainty_mb
\end{verbatim}
The energy grid must resolve reaction thresholds and the structure below
\SI{50}{MeV}, which has the strongest leverage on the distal activity fall-off.
Each table will carry its literature or database provenance and interpolation
convention.

\subsection{Public cross-section inputs}

The alternative production curves will be assembled from three public nuclear
data sources.  The imported curves form the numerator of
Eq.~\eqref{eq:weight}; the Geant4 denominator is measured with the thin-target
runs described below.

\subsubsection*{EXFOR measurements}
The IAEA Nuclear Data Section publishes the experimental EXFOR library in
machine-readable form.  The most direct representation for this work is the
\href{https://github.com/IAEA-NDS/EXFOR-Plot-File}{EXFOR Plot File repository}.
Its Compact EXFOR (CX4) files contain the incident energy, energy uncertainty,
cross section and cross-section uncertainty, while the directory and file names
retain the projectile, target, product, author, year and EXFOR subentry.  The
initial selection is
\begin{verbatim}
cx4/p/C12/x/sig/*prodC11*.cx4
cx4/p/O16/x/sig/*prodO15*.cx4
cx4/p/O16/x/sig/*prodC11*.cx4
cx4/p/N14/x/sig/*prodN13*.cx4   # optional N-13 extension
\end{verbatim}
The corresponding EXFOR entry metadata are available from the
\href{https://github.com/IAEA-NDS/exfor_master}{EXFOR Master Files repository},
which is distributed under the CC BY 4.0 licence.

Each experimental series will remain a separate versioned input, identified by
its EXFOR entry and subentry.  The labels ``EXFOR 1'' and ``EXFOR 2'' in
Espa\~na et al. denote curves assembled from selected experiments over different
energy intervals rather than released database products.  A reproduction of
either curve therefore requires an explicit record of the selected EXFOR
subentries, interpolation and joining energies.  The primary uncertainty
ensemble will use the individual EXFOR series, including measurements published
after that study.

\subsubsection*{TENDL evaluated curves}
\href{https://tendl.imperial.ac.uk/tendl_2023/tendl2023.html}{TENDL-2023}
provides TALYS-based evaluated data up to \SI{200}{MeV}, including covariance
information.  The relevant incident-proton target files are
\href{https://tendl.imperial.ac.uk/tendl_2023/proton_html/C/ProtonC12.html}
{C-12},
\href{https://tendl.imperial.ac.uk/tendl_2023/proton_html/O/ProtonO16.html}
{O-16} and
\href{https://tendl.imperial.ac.uk/tendl_2023/proton_html/N/ProtonN14.html}
{N-14}.  Residual-production curves for C-11, O-15 and, when included, N-13
will be extracted from these evaluated ENDF files.

\subsubsection*{JENDL evaluated curves}
The
\href{https://wwwndc.jaea.go.jp/ftpnd/jendl/jendl40he_prod.html}
{JENDL-4.0/HE residual-production collection} supplies precomputed production
cross sections for residual nuclei.  It constructs these tables from the
reaction cross section and the relative residual-product yield stored in the
evaluated file.  Direct proton-target tables are publicly available for
\href{https://wwwndc.jaea.go.jp/ftpnd/jendl/jendl40he/out_p/C012.txt}{C-12},
\href{https://wwwndc.jaea.go.jp/ftpnd/jendl/jendl40he/out_p/O016.txt}{O-16} and
\href{https://wwwndc.jaea.go.jp/ftpnd/jendl/jendl40he/out_p/N014.txt}{N-14}.

ICRU Report 63 remains a useful historical comparator because it was used in
the Espa\~na study.  Its official distribution is attached to the purchased
report.  The reproducible public input ensemble is therefore EXFOR, TENDL-2023
and JENDL-4.0/HE; an ICRU curve can be added when its licensed numerical files
are available.

For every imported table, the local provenance record will contain the source
library and release, retrieval date, source URL, dataset or subentry identifier,
original reaction code, author and year, original units, independent or
cumulative status, interpolation rule and any transformation applied during
conversion to the common CSV schema.  Experimental measurements will retain
their quoted energy and cross-section uncertainties.

Geant4 exposes the total inelastic interaction cross section, while a
residual-specific production cross section is generated by the interaction
model.  Effective Geant4 channel curves will therefore be measured with
thin-target simulations at fixed proton energies.  Carbon and oxygen targets
provide the three dominant channels in Eq.~\eqref{eq:channels}; nitrogen can be
added for the $\mathrm{^{13}N}$ component.  The calculation will use the same
Geant4 version and nominal physics list as the treatment simulation.

\subsection{Energy grid}

The initial scan uses the common 30-point grid in
Table~\ref{tab:energy-grid}.  It extends below the production thresholds and
beyond the maximum energy of the current SOBP.  The fine spacing between
\SI{12}{MeV} and \SI{30}{MeV} resolves the channel turn-on, while
\SI{5}{MeV} spacing covers the distal-edge and treatment-energy regions.

\begin{table}[htbp]
  \centering
  \caption{Proposed proton-energy grid for the thin-target runs.}
  \label{tab:energy-grid}
  \begin{tabular}{
    >{\raggedright\arraybackslash}p{0.22\textwidth}
    >{\raggedright\arraybackslash}p{0.48\textwidth}
    >{\raggedright\arraybackslash}p{0.20\textwidth}}
    \toprule
    Region & Proton energies (MeV) & Purpose \\
    \midrule
    Below threshold
      & 5, 10
      & Confirm zero production \\
    Threshold region
      & 12, 14, 16, 18, 20, 22, 24, 26, 28, 30
      & Resolve channel turn-on \\
    Distal-edge region
      & 35, 40, 45, 50, 55, 60
      & Dense low-energy coverage \\
    Treatment range
      & 65, 70, 75, 80, 85, 90, 95, 100, 105, 110, 115, 120
      & Cover the SOBP with margin \\
    \bottomrule
  \end{tabular}
\end{table}

After the first pass, additional \SI{1}{MeV} points will be inserted around a
sharply resolved threshold or peak.  The cross-section table will use the
actual proton kinetic energy immediately before the producing interaction,
rather than the nominal energy at the target entrance.

\subsection{Target thickness}

The baseline target areal thickness is
\begin{equation}
  \tau = \SI{5}{mg/cm^2}.
  \label{eq:areal-thickness}
\end{equation}
For material density $\rho$, the physical thickness is
\begin{equation}
  \ell = \frac{\tau}{\rho}.
  \label{eq:physical-thickness}
\end{equation}
Table~\ref{tab:targets} gives representative physical thicknesses and nuclear
areal densities.  The Geant4 material definition will state the density and
isotopic composition used in each run.

\begin{table}[htbp]
  \centering
  \caption{Baseline thin targets for an areal thickness of
  \SI{5}{mg/cm^2}.}
  \label{tab:targets}
  \begin{tabular}{lccc}
    \toprule
    Target & Density (\si{g/cm^3}) & Thickness ($\mu$m) &
      Nuclear areal density (\si{cm^{-2}}) \\
    \midrule
    $\Cl{12}$ graphite & 2.00 & 25 & \num{2.51e20} \\
    $\Ox{16}$          & 1.14 & 44 & \num{1.88e20} \\
    $\mathrm{^{14}N}$  & 0.81 & 62 & \num{2.15e20} \\
    \bottomrule
  \end{tabular}
\end{table}

In the thin-target limit, the effective Geant4 channel cross section is
\begin{equation}
  \sigma_c^{(\mathrm{G4})}(E) \simeq
  \frac{N_c(E)}{N_p(E)\,n_t\,\ell},
  \label{eq:thin-target-xs}
\end{equation}
where $N_c$ is the number of residual nuclei produced through channel $c$,
$N_p$ is the number of incident protons, and $n_t\ell$ is the target nuclear
areal density.  For a representative cross section of \SI{500}{mb}, the targets
in Table~\ref{tab:targets} give an interaction probability of order
\begin{equation}
  p_{\mathrm{int}} \simeq n_t\ell\,\sigma_{\mathrm{inel}}
  \sim 10^{-4},
\end{equation}
so multiple nuclear interactions are negligible.

The target design is accepted when
\begin{equation}
  p_{\mathrm{inel}} < 10^{-3},
  \qquad
  \langle\Delta E\rangle < \SI{0.5}{MeV}.
  \label{eq:target-acceptance}
\end{equation}
If the low-energy samples exceed the energy-loss criterion, the areal thickness
below \SI{30}{MeV} will be reduced to \SI{2}{mg/cm^2}; the
\SI{5}{mg/cm^2} target will be retained at and above \SI{30}{MeV}.

\section{Cross-section-only reweighting}

Let emitter $i$ be produced through channel $c_i$ at projectile energy $E_i$.
For cross-section model $m$, assign the weight
\begin{equation}
  w_i^{(m)} =
  \frac{\sigma_{c_i}^{(m)}(E_i)}
       {\sigma_{c_i}^{(\mathrm{G4})}(E_i)}.
  \label{eq:weight}
\end{equation}
The model-dependent production of isotope $j$, normalized to a delivered dose
$D$, is then
\begin{equation}
  P_j^{(m)}(D) =
  \frac{D}{D_{\mathrm{sim}}}
  \sum_{i\in j}w_i^{(m)},
  \label{eq:production}
\end{equation}
where $D_{\mathrm{sim}}$ is the nominal target dose scored by Stage A.

The spatial distribution for isotope $j$ is sampled with probability
\begin{equation}
  p_i^{(m)} =
  \frac{w_i^{(m)}}{\sum_{k\in j}w_k^{(m)}}.
  \label{eq:sampling}
\end{equation}
Equations~\eqref{eq:production} and \eqref{eq:sampling} modify both the isotope
budget and the depth-dependent source shape.  Radioactive decay, acquisition
timing, positron transport, detector response and reconstruction retain their
existing treatment.

A Python tool in \file{analysis\_transport/} will read the emitter provenance
and cross-section tables and write
\begin{itemize}[nosep]
  \item one per-emitter weight vector for each model;
  \item weighted production summaries by isotope and reaction channel;
  \item weighted acquisition budgets for the requested timing scenarios;
  \item nominal and alternative depth profiles with their distal-edge shifts.
\end{itemize}
The downstream source sampler will use Eq.~\eqref{eq:sampling}.  Common random
numbers will pair nominal and alternative acquisitions, reducing the Monte Carlo
noise in Eq.~\eqref{eq:bias}.

\section{Validation}

The following checks define acceptance of the implementation:
\begin{enumerate}
  \item Every listed emitter has a valid residual-production record, apart from
        explicitly classified production routes.
  \item Setting all weights to unity reproduces the current isotope counts,
        budgets and spatial profiles.
  \item Thin-target yields reproduce the effective Geant4 channel curves used
        in the denominator of Eq.~\eqref{eq:weight}.
  \item Applying a known analytic weight reproduces the expected weighted yield
        and energy dependence.
  \item The nominal depth-dose distribution and $D_{\mathrm{sim}}$ are invariant
        under emitter weighting.
  \item Alternative physics lists write separate run directories, and the run
        metadata records the selected list.
  \item Existing scenario readers continue to read the established emitter
        coordinates and isotope identifier.
  \item Repeated runs with fixed seeds reproduce the weighted budgets and paired
        distal-edge shifts.
\end{enumerate}

\section{Execution plan}

\begin{enumerate}
  \item Implement production provenance and run-time physics-list selection.
  \item Extend \file{common/SCHEMA.md}, the scenario schema and the snapshot
        metadata to describe the new fields.
  \item Run a \num{1e6}-proton \file{uniform\_headep} pilot and validate
        provenance completeness, channel classification and backward-compatible
        output.
  \item Produce the thin-target Geant4 channel curves and import the alternative
        cross-section data.
  \item Implement the weight and weighted-budget calculation, including the
        unity-weight regression test.
  \item Propagate pilot samples through the existing acquisition and distal-edge
        estimator.
  \item Generate the new \num{1e8}-proton nominal
        \file{uniform\_headep\_sobp} source with production provenance.
  \item Generate the BIC, INCL++ and Bertini comparison runs and compute the
        hadronic-model envelope.
  \item Report the individual $b_{\mathrm{xs}}^{(m)}$, the adopted cross-section
        uncertainty and the broader hadronic-model envelope.
\end{enumerate}

\section{Principal code changes}

The implementation is concentrated in the following files:
\begin{itemize}[nosep]
  \item \file{stageA\_transport/proton\_transport.cc}: physics-list selection;
  \item \file{StageAConfig.hh}: nominal physics-list name;
  \item \file{ActionInitialization.cc}: connect event and stepping provenance;
  \item \file{EventAction.hh}: residual-track production records;
  \item \file{SteppingAction.cc}: hadronic target, projectile and secondary
        capture;
  \item \file{TrackingAction.cc}: association with positron annihilation;
  \item \file{EmitterData.hh}: provenance columns and MT merge;
  \item \file{RunAction.cc}: CSV output, dynamic metadata and collision-free run
        tags;
  \item \file{decay\_sampling/budget.py}: weighted isotope production;
  \item \file{analysis\_transport/}: thin-target analysis, cross-section weights
        and validation;
  \item \file{common/SCHEMA.md} and \file{tools/scenario\_template/SCHEMA.md}:
        shared data contract.
\end{itemize}

\end{document}
